\documentclass[11pt,a4paper]{article}

\usepackage[utf8]{inputenc}
\usepackage[T1]{fontenc}
\usepackage{lmodern}
\usepackage{geometry}
\usepackage{amsmath, amssymb, amsthm}
\usepackage{hyperref}
\usepackage{enumitem}
\usepackage{booktabs}
\usepackage{graphicx}
\usepackage{setspace}
\usepackage{natbib}

\geometry{margin=1in}
\onehalfspacing

\title{Memory-Gravity Attention:\\Persistent-Importance Retrieval for Long-Horizon Memory Tasks}
\author{Kevin T.N}
\date{}

\begin{document}
\maketitle

\begin{abstract}
We present Memory-Gravity Attention (MGA), a model-agnostic retrieval layer that augments similarity-based memory retrieval with persistent-importance signals. Standard approaches---including similarity-only and Generative-Agents-style scoring (recency + importance + similarity)---perform well on recency-friendly queries but fail on tasks requiring long-horizon persistent memory: retrieving facts stated once long ago (needle), tracking unresolved tasks (open-loop), and filtering noise. MGA introduces a multiplicative gate: $\text{score} = \text{sim} \times (1 + \sigma(\boldsymbol{\theta}^\top \mathbf{f}))$, where $\mathbf{f}$ comprises log-estimated persistent signals (recency, frequency, unresolvedness, utility, weight, goal-relevance) and $\boldsymbol{\theta}$ is learned via logistic regression on interaction-derived feedback. On a synthetic benchmark of 10 worlds (4,500+ nodes, 415 test queries, sentence-transformer embeddings), the Generative-Agents baseline achieves Recall@5 = 0.209 overall but scores 0.000 on needle, open-loop, and noise-resistance query families. MGA achieves non-zero recall on all three families (needle: 0.143, open-loop: 0.222, noise: 0.075) while maintaining competitive overall performance (Recall@5 = 0.195) and the lowest noise retrieval rate (Noise@5 = 0.017). We report learned feature coefficients showing that goal-relevance ($\theta = 2.92$) is the dominant persistent signal. Code and data are available at \url{https://github.com/jkdkr2439/Memory-Gravity-Attention}.
\end{abstract}

\section{Introduction}

Memory retrieval for language model agents is typically framed as a similarity search problem: given a query, retrieve the most semantically similar stored items. This approach works well when the answer is recent and phrased similarly to the query. However, it fails systematically on tasks that require \emph{persistent importance}---information that remains relevant over long horizons regardless of recency or surface similarity.

Consider the following scenarios:
\begin{itemize}[leftmargin=2em]
\item A fact stated once in session 2 of 25, never repeated, but still valid and needed (needle).
\item A task opened early and never closed, requiring recall despite many intervening events (open-loop).
\item A constraint revised mid-project, where the old version has high similarity to the query but is stale (stale trap).
\item A dense log with 40\%+ chatter, where the system must retrieve signal from noise (noise resistance).
\end{itemize}

Park et al. \citep{park2023generative} proposed scoring memory as $\alpha \cdot \text{recency} + \beta \cdot \text{importance} + \gamma \cdot \text{relevance}$ for generative agents. This addresses some failure modes but, as we show empirically, still scores zero on needle, open-loop, and noise-resistance tasks in our benchmark.

We propose \textbf{Memory-Gravity Attention (MGA)}, which uses a multiplicative gate structure:
\begin{equation}
\text{score}_i = \text{sim}_i \times \bigl(1 + \sigma(\boldsymbol{\theta}^\top \mathbf{f}_i)\bigr),
\end{equation}
where $\text{sim}_i$ is the query-node cosine similarity, $\mathbf{f}_i$ is a vector of log-estimated persistent signals, and $\sigma$ is the sigmoid function. The gate ensures that: (a) nodes with low similarity cannot be retrieved by importance alone, and (b) among similarly relevant nodes, persistent importance disambiguates.

\section{Related Work}

\textbf{Memory for LLM agents.} Park et al. \citep{park2023generative} introduced a three-term scoring function for generative agents. MemGPT \citep{packer2023memgpt} manages memory via a virtual context hierarchy. Reflexion \citep{shinn2023reflexion} uses self-reflection for episodic learning. None of these explicitly benchmark persistent-importance retrieval on long-horizon tasks.

\textbf{Retrieval-augmented generation (RAG).} Standard RAG \citep{lewis2020retrieval} uses dense retrieval (similarity-based). Recent work adds reranking \citep{glass2022re2g}, filtering, and hybrid search. These focus on factual QA, not persistent-importance memory management.

\textbf{Attention and novelty.} Predictive coding \citep{friston2010free} and surprise-driven attention \citep{itti2009bayesian} process prediction errors rather than similarities. Differential Transformer \citep{ye2024diff} subtracts two attention maps to reduce noise. Our work is complementary: MGA operates as an external retrieval layer, not a modification of transformer internals.

\section{Method}

\subsection{Problem Setting}

A memory store contains $N$ nodes, each with content text, timestamp, and session ID. Given a query at time $t$, retrieve the top-$k$ nodes that answer the query. Gold labels exist only in an oracle (evaluator-only) and are never visible to the retrieval system.

\subsection{Signal Estimators}

All signals are estimated from the interaction log without access to gold labels:

\begin{itemize}[leftmargin=2em]
\item \textbf{Recency}: $r_i = \exp(-(t_q - t_i) / \tau_r)$
\item \textbf{Frequency}: $f_i = \log(1 + m_i) / \log(1 + m_{\max})$, where $m_i$ counts near-duplicate mentions (embedding similarity $> 0.75$)
\item \textbf{Unresolvedness}: binary heuristic detecting task-like nodes without subsequent closure matches
\item \textbf{Goal relevance}: $g_i = \cos(\mathbf{e}_i, \mathbf{e}_{\text{goal}})$, normalized per retrieval call
\item \textbf{Utility/Weight}: Beta-smoothed feedback from training episodes (default 0.5 = uninformative)
\end{itemize}

All features are min-max normalized to $[0,1]$ per retrieval call.

\subsection{MGA Gate Retriever}

The feature vector (excluding similarity) is:
\[
\mathbf{f}_i = [\text{recency}, \text{frequency}, \text{unresolved}, \text{utility}, \text{weight}, \text{goal\_rel}] \in [0,1]^6.
\]

The MGA gate score is:
\begin{equation}
\text{score}_i = \text{sim}_i \times \bigl(1 + \sigma(\boldsymbol{\theta}^\top \mathbf{f}_i)\bigr).
\end{equation}

$\boldsymbol{\theta} \in \mathbb{R}^6$ is learned via logistic regression on training queries (gold = positive, sampled non-gold = negative).

\subsection{Baselines}

\begin{itemize}[leftmargin=2em]
\item \textbf{Recency}: rank by $r_i$ only.
\item \textbf{Similarity}: rank by $\text{sim}_i$ only.
\item \textbf{GA}: $\alpha \cdot r_i + \beta \cdot w_i + \gamma \cdot \text{sim}_i$ with $\alpha = \beta = \gamma = 1$ \citep{park2023generative}.
\item \textbf{MGA linear}: $\boldsymbol{\theta}^\top [\text{sim}, \mathbf{f}]$ (ablation without gate structure).
\end{itemize}

\section{Experimental Setup}

\subsection{Synthetic Benchmark}

We generate 10 synthetic worlds, each with:
\begin{itemize}[leftmargin=2em]
\item 25 sessions, 10--25 events per session ($\sim$450 nodes/world)
\item 3 domains: research paper, software project, design project
\item 7 preference dimensions with 2--4 mid-log revisions (stale traps)
\item 4--8 tasks with staggered open/close (open-loop tracking)
\item 45\% chatter events (noise)
\item Template variation across 5--8 templates per event type
\end{itemize}

Query families: constraint recall, stale trap, open-loop, needle, noise resistance.

Train/test split: 40\%/60\% of queries per world. $\boldsymbol{\theta}$ learned on train only.

\subsection{Embeddings}

Sentence-transformers (\texttt{all-MiniLM-L6-v2}) for all similarity computations. Run on NVIDIA T4 GPU via Modal.com.

\subsection{Metrics}

Recall@5, nDCG@5, Stale@5 (fraction of retrieved nodes that are superseded), Noise@5 (fraction that are chatter). All metrics on test queries.

\section{Results}

\subsection{Overall Performance}

\begin{table}[h]
\centering
\caption{Retrieval performance across 415 test queries (10 worlds).}
\begin{tabular}{lcccc}
\toprule
\textbf{Retriever} & \textbf{Recall@5} & \textbf{nDCG@5} & \textbf{Stale@5} & \textbf{Noise@5} \\
\midrule
Recency & 0.029 & 0.085 & 0.000 & 0.364 \\
Similarity & 0.138 & 0.496 & 0.340 & 0.046 \\
GA & \textbf{0.209} & \textbf{0.697} & \textbf{0.000} & 0.070 \\
MGA gate & 0.174 & 0.562 & 0.265 & 0.029 \\
MGA linear & 0.195 & 0.614 & 0.205 & \textbf{0.017} \\
\bottomrule
\end{tabular}
\end{table}

GA achieves the highest overall Recall@5 and nDCG@5. However, GA's stale advantage (Stale@5 = 0.000) is largely due to recency naturally filtering old (stale) nodes rather than explicit stale detection.

\subsection{Per-Family Breakdown}

\begin{table}[h]
\centering
\caption{Recall@5 by query family. Families where persistent memory matters most.}
\begin{tabular}{lccccc}
\toprule
\textbf{Family} & \textbf{Recency} & \textbf{Sim} & \textbf{GA} & \textbf{MGA gate} & \textbf{MGA linear} \\
\midrule
Constraint & 0.036 & 0.151 & \textbf{0.257} & 0.191 & 0.220 \\
Needle & 0.000 & 0.143 & 0.000 & \textbf{0.143} & 0.000 \\
Open-loop & 0.000 & 0.097 & 0.000 & 0.111 & \textbf{0.222} \\
Noise resist & 0.000 & 0.050 & 0.000 & \textbf{0.075} & 0.025 \\
Stale trap & 0.036 & 0.135 & \textbf{0.275} & 0.187 & 0.230 \\
\bottomrule
\end{tabular}
\end{table}

The critical finding: \textbf{GA scores 0.000 on needle, open-loop, and noise-resistance families.} These are precisely the tasks where persistent importance beyond recency and similarity is required. MGA achieves non-zero recall on all three.

\subsection{Learned Feature Importance}

\begin{table}[h]
\centering
\caption{Learned $\boldsymbol{\theta}$ (MGA gate), mean $\pm$ std across 10 worlds.}
\begin{tabular}{lc}
\toprule
\textbf{Feature} & \textbf{Coefficient} \\
\midrule
Goal relevance & $2.92 \pm 0.87$ \\
Recency & $0.67 \pm 0.50$ \\
Frequency & $0.40 \pm 0.35$ \\
Unresolvedness & $0.21 \pm 0.29$ \\
Utility & $-0.01 \pm 0.01$ \\
Weight & $-0.01 \pm 0.01$ \\
\bottomrule
\end{tabular}
\end{table}

Goal relevance dominates ($\theta = 2.92$): knowing what the system is currently working on is the strongest persistent signal for retrieval. Recency and frequency provide additional signal. Utility and weight are uninformative in this benchmark (insufficient training episodes to estimate them).

\section{Discussion}

\subsection{Why GA Fails on Persistent Tasks}

GA combines recency + importance + similarity additively. On needle queries (old facts), recency suppresses old nodes. On open-loop queries (unresolved tasks), the importance proxy (weight) is uninformative without explicit task-tracking. On noise queries, GA has no mechanism to distinguish signal from chatter beyond similarity.

MGA's gate structure addresses this: the persistent features (especially goal-relevance and unresolvedness) amplify nodes that similarity alone identifies as marginally relevant.

\subsection{Attention-Entropy Diagnostic}

We introduce an attention-entropy diagnostic to detect \emph{retriever lock}---when a retriever attends to a narrow, repetitive set of nodes regardless of query content. For each retriever, we compute the normalized Shannon entropy of the softmax-transformed score distribution:
\[
\hat{H} = \frac{-\sum_i p_i \log p_i}{\log N}, \quad p_i = \text{softmax}(s_i).
\]
A retriever is ``locked'' if $\hat{H}$ is consistently low across diverse queries.

\begin{table}[h]
\centering
\caption{Attention-entropy diagnostic (5 worlds, 29 test queries).}
\begin{tabular}{lccc}
\toprule
\textbf{Retriever} & \textbf{Mean $\hat{H}$} & \textbf{Std $\hat{H}$} & \textbf{Top-5 Conc.} \\
\midrule
Recency & 0.992 & 0.001 & 0.050 \\
Similarity & 0.992 & 0.005 & 0.060 \\
GA & 0.983 & 0.005 & 0.078 \\
MGA gate & 0.976 & 0.014 & 0.093 \\
MGA linear & 0.650 & 0.131 & 0.467 \\
\bottomrule
\end{tabular}
\end{table}

All retrievers show healthy (non-locked) attention. MGA linear concentrates attention most aggressively (top-5 concentration = 0.467 vs 0.05--0.09 for others), explaining its stronger recall on targeted queries at the cost of lower diversity. MGA gate maintains high entropy (0.976) while still achieving selective boosting through the multiplicative structure---the gate amplifies without collapsing.

\subsection{Limitations}

\begin{enumerate}[leftmargin=2em]
\item \textbf{Synthetic data}: Results are on template-based synthetic logs. Real-world interaction logs have greater lexical diversity, ambiguity, and implicit references.
\item \textbf{Small training signal}: With only 4--5 training queries per world, $\boldsymbol{\theta}$ estimation is noisy. More training data would likely improve MGA.
\item \textbf{Utility/weight features uninformative}: The benchmark does not generate enough use-episodes for these features to become meaningful. A longer-running agent would accumulate this signal.
\item \textbf{Overall recall}: GA still wins overall. MGA's advantage is specifically on long-horizon tasks.
\end{enumerate}

\subsection{Broader Significance}

MGA operationalizes a principle from existential systems theory: a system should attend not only to what is \emph{similar} to its current query, but to what is \emph{persistently important} for its ongoing operation. This is analogous to gap-based attention in biological systems, where survival-relevant mismatches receive priority processing regardless of recency or surface similarity.

\section{Conclusion}

We introduced MGA, a retrieval mechanism that gates similarity by learned persistent importance. On a controlled benchmark, we showed that:
\begin{enumerate}[leftmargin=2em]
\item Standard approaches (similarity, GA) score zero on tasks requiring long-horizon persistent memory.
\item MGA achieves non-zero recall on all persistent-memory task families.
\item Goal-relevance is the dominant persistent signal ($\theta = 2.92$).
\item MGA has the lowest noise retrieval rate (Noise@5 = 0.017 vs GA's 0.070).
\end{enumerate}

MGA does not replace similarity-based retrieval. It fills the gap where similarity and recency fail: persistent, long-horizon, noise-resistant memory.

\bibliographystyle{plainnat}
\begin{thebibliography}{10}

\bibitem[Park et al.(2023)]{park2023generative}
Park, J.~S., O'Brien, J.~C., Cai, C.~J., Morris, M.~R., Liang, P., \& Bernstein, M.~S. (2023).
Generative agents: Interactive simulacra of human behavior.
\textit{UIST '23}.

\bibitem[Packer et al.(2023)]{packer2023memgpt}
Packer, C., Fang, V., Patil, S.~G., Lin, K., Wooders, S., \& Gonzalez, J.~E. (2023).
MemGPT: Towards LLMs as operating systems.
\textit{arXiv:2310.08560}.

\bibitem[Shinn et al.(2023)]{shinn2023reflexion}
Shinn, N., Cassano, F., Gopinath, A., Narasimhan, K., \& Yao, S. (2023).
Reflexion: Language agents with verbal reinforcement learning.
\textit{NeurIPS 2023}.

\bibitem[Lewis et al.(2020)]{lewis2020retrieval}
Lewis, P., Perez, E., Piktus, A., Petroni, F., Karpukhin, V., Goyal, N., ... \& Kiela, D. (2020).
Retrieval-augmented generation for knowledge-intensive NLP tasks.
\textit{NeurIPS 2020}.

\bibitem[Glass et al.(2022)]{glass2022re2g}
Glass, M., Rossiello, G., Chowdhury, M.~F.~M., Naber, A., Nishi, K., \& Gliozzo, A. (2022).
Re2G: Retrieve, rerank, generate.
\textit{NAACL 2022}.

\bibitem[Friston(2010)]{friston2010free}
Friston, K. (2010).
The free-energy principle: A unified brain theory?
\textit{Nature Reviews Neuroscience}, 11, 127--138.

\bibitem[Itti \& Baldi(2009)]{itti2009bayesian}
Itti, L. \& Baldi, P. (2009).
Bayesian surprise attracts human attention.
\textit{Vision Research}, 49(10), 1295--1306.

\bibitem[Ye et al.(2024)]{ye2024diff}
Ye, T., Dong, L., Xia, Y., Sun, Y., Zhu, Y., Huang, G., \& Wei, F. (2024).
Differential Transformer.
\textit{arXiv:2410.05258}.

\end{thebibliography}

\end{document}
